% Presentation: CNN Performance Analysis (Final Results)
\documentclass{beamer}

\usepackage[T1]{fontenc}
\usepackage[brazil]{babel}
\usepackage[utf8]{inputenc}
\usepackage{booktabs}
\usepackage{graphicx}

% Choose the Inf theme
\usetheme{Inf}

% Title Information
\title[Performance Analysis of CNNs]
      {Performance Analysis of Convolutional Neural Network Architectures for Image Classification}

\subtitle{Comp. Sys. Perf. Analysis}

\date{December, 2025}

\author{Arthur Dohnert, Eduardo Menezes, \\ Fernando Tedesco}
\institute{Instituto de Informática --- UFRGS}

\begin{document}

% 1. Title Slide
\InfTitlePage

% 2. Agenda
\begin{frame}
  \frametitle{Content}
  \tableofcontents
\end{frame}

% =========================================================
% SECTION 1: Introduction & Objectives
% =========================================================
\section{Introduction}

\begin{frame}
    \frametitle{Context \& Motivation}
    \begin{itemize}
        \item \textbf{Problem}: Deploying CNNs in resource-constrained environments (Edge, Embedded, Shared GPU Clusters).
        \item \textbf{Trade-off}: High predictive performance vs. low latency and memory usage.
        \item \textbf{Goal}: Systematically benchmark lightweight architectures under a unified, reproducible pipeline.
    \end{itemize}
\end{frame}

\begin{frame}
    \frametitle{Evaluated Architectures}
    \begin{columns}
        \begin{column}{0.5\textwidth}
            \textbf{Light Models (Green)}
            \begin{itemize}
                \item MobileNet V1
                \item ShuffleNet V2
                \item SqueezeNet
                \item EfficientNet-B0
            \end{itemize}
            \vspace{0.4cm}
            \textbf{Medium Models (Blue)}
            \begin{itemize}
                \item VGG16
                \item ResNet34
                \item Inception V3
                \item DenseNet121
            \end{itemize}
        \end{column}

        \begin{column}{0.5\textwidth}
            \textbf{Heavy Models (Red)}
            \begin{itemize}
                \item ResNet101
                \item Inception V4
                \item Xception
                \item EfficientNet-B7
            \end{itemize}
            
        \end{column}
    \end{columns}
\end{frame}

% =========================================================
% SECTION 2: Methodology
% =========================================================
\section{Methodology}

\begin{frame}
    \frametitle{Experimental Setup}
    \begin{itemize}
        \item \textbf{Hardware}: UFRGS PCAD Cluster (Tupi partition).
            \begin{itemize}
                \item Single NVIDIA GPU (24 GiB VRAM) per job.
                \item Data staged to node-local \texttt{SCRATCH} to minimize I/O latency.
            \end{itemize}
        \item \textbf{Software}:
            \begin{itemize}
                \item Python 3.11, PyTorch 2.0.
                \item \texttt{PYTORCH\_CUDA\_ALLOC\_CONF} tuned to reduce fragmentation.
            \end{itemize}
        \item \textbf{Dataset}: Tiny ImageNet (200 classes, $64\times64$ px).
        \item \textbf{Protocol}: 100 Epochs, 5 Random Seeds per model.
    \end{itemize}
\end{frame}


% =========================================================
% SECTION 3: Results
% =========================================================
\section{Results} % faltam muitos gráficos. vamos botar uma página para cada gráfico. não sei se aquela tabela é necessária, se sim, vamos arrumar os dados, até porque agora temos 12 modelos

\subsection{Validation Results}

\begin{frame}
    \frametitle{Validation Accuracy vs. Epoch}
    \begin{figure}
        \centering
        % Ensure fig_val_accuracy.png is in the folder
        \includegraphics[width=1\textwidth]{../reports/results/accuracy_vs_epoch.png}
    \end{figure}
\end{frame}

\begin{frame}
    \frametitle{F1 Score vs. Epoch}
    \begin{figure}
        \centering
        % Ensure fig_val_accuracy.png is in the folder
        \includegraphics[width=1\textwidth]{../reports/results/f1_vs_epoch.png}
    \end{figure}
\end{frame}

\begin{frame}
    \frametitle{Peak Performance Summary}
    \begin{table}
        \centering
        \begin{tabular}{lcc}
            \toprule
            \textbf{Model} & \textbf{Top-1 Accuracy} & \textbf{F1-Score} \\
            \midrule
            DenseNet121      & $0.184 \pm 0.015$ & $0.154 \pm 0.015$ \\
            EfficientNet-B0  & $0.603 \pm 0.002$ & $0.598 \pm 0.002$ \\
            EfficientNet-B7  & $0.463 \pm 0.017$ & $0.449 \pm 0.018$ \\
            Inception V3     & $0.670 \pm 0.002$ & $0.669 \pm 0.002$ \\
            Inception V4     & $0.557 \pm 0.011$ & $0.553 \pm 0.011$ \\
            MobileNet V1     & $0.606 \pm 0.002$ & $0.604 \pm 0.002$ \\
            ResNet101        & $0.207 \pm 0.102$ & $0.181 \pm 0.092$ \\
            ResNet34         & $0.633 \pm 0.005$ & $0.630 \pm 0.005$ \\
            ShuffleNet V2    & $0.506 \pm 0.008$ & $0.498 \pm 0.008$ \\
            SqueezeNet       & $0.147 \pm 0.071$ & $0.123 \pm 0.062$ \\
            VGG16            & $0.567 \pm 0.004$ & $0.561 \pm 0.004$ \\
            Xception         & $0.669 \pm 0.003$ & $0.667 \pm 0.003$ \\
            \bottomrule
        \end{tabular}
        \caption{Peak metrics (Mean $\pm$ 95\% CI) across 5 seeds.}
    \end{table}
\end{frame}

\subsection{GPU Usage}

\begin{frame}
    \frametitle{GPU Utilization vs. Epoch}
    \begin{figure}
        \centering
        % Ensure fig_val_accuracy.png is in the folder
        \includegraphics[width=1\textwidth]{../reports/results/gpu_utilization.png}
    \end{figure}
\end{frame}
\begin{frame}
    \frametitle{GPU VRAM Usage vs. Epoch}
    \begin{figure}
        \centering
        % Ensure fig_val_accuracy.png is in the folder
        \includegraphics[width=1\textwidth]{../reports/results/gpu_vram_usage.png}
    \end{figure}
\end{frame}
\begin{frame}
    \frametitle{GPU Temperature vs. Epoch}
    \begin{figure}
        \centering
        % Ensure fig_val_accuracy.png is in the folder
        \includegraphics[width=1\textwidth]{../reports/results/gpu_temperature.png}
    \end{figure}
\end{frame}

\subsection{Time Consumption}

\begin{frame}
    \frametitle{Total Training Time}
    \begin{figure}
        \centering
        % Ensure fig_val_accuracy.png is in the folder
        \includegraphics[width=1\textwidth]{../reports/results/training_time_total_bar.png}
    \end{figure}
\end{frame}
\begin{frame}
    \frametitle{Total Energy Consumption per Model}
    \begin{figure}
        \centering
        % Ensure fig_val_accuracy.png is in the folder
        \includegraphics[width=1\textwidth]{../reports/results/total_energy_per_model_weighted.png}
    \end{figure}
\end{frame} 

\begin{frame}
    \frametitle{Training Time Distribution}
    \begin{figure}
        \centering
        % Ensure fig_val_accuracy.png is in the folder
        \includegraphics[width=1\textwidth]{../reports/results/training_time_boxplot.png}
    \end{figure}
\end{frame}
\begin{frame}
    \frametitle{Training Time per Seed}
    \begin{figure}
        \centering
        % Ensure fig_val_accuracy.png is in the folder
        \includegraphics[width=1\textwidth]{../reports/results/training_time_per_seed.png}
    \end{figure}
\end{frame}
\begin{frame}
    \frametitle{Average Inference Time per Model}
    \begin{figure}
        \centering
        % Ensure fig_val_accuracy.png is in the folder
        \includegraphics[width=1\textwidth]{../reports/results/inference_time_per_model.png}
    \end{figure}
\end{frame}  

\subsection{Trade - Offs}
\begin{frame}
    \frametitle{Accuracy vs. Training Cost}
    \begin{figure}
        \centering
        % Ensure fig_val_accuracy.png is in the folder
        \includegraphics[width=1\textwidth]{../reports/results/accuracy_vs_training_cost.png}
    \end{figure}
\end{frame}
\begin{frame}
    \frametitle{Accuracy vs. Energy Consumption}
    \begin{figure}
        \centering
        % Ensure fig_val_accuracy.png is in the folder
        \includegraphics[width=1\textwidth]{../reports/results/accuracy_vs_total_energy_weighted.png}
    \end{figure}
\end{frame}
% =========================================================
% SECTION 4: Conclusion
% =========================================================
\section{Conclusion}

\begin{frame}
    \frametitle{Key Takeaways}
    \begin{enumerate}
        \item \textbf{Performance does not strictly follow model size.}  
              Surprisingly, several heavy models underperformed compared to lightweight architectures.  
              Inception V3, ResNet34, MobileNet V1, and Xception achieved the best accuracies, despite having very different computational footprints.

        \item \textbf{Training time was also inconsistent with model complexity.}  
              Some lightweight models (e.g., SqueezeNet) trained slower than expected,  
              while certain heavier models (e.g., Inception V3, Xception) trained efficiently.
    \end{enumerate}
\end{frame}


\begin{frame}
    \frametitle{Study Limitations}
    \begin{itemize}
        \item \textbf{Fixed Hyperparameters.}  
              All models were trained using identical hyperparameters.  
              This design isolates architectural effects but likely  
              underestimates the true performance of models that  
              require tailored tuning (e.g., SqueezeNet, EfficientNet).

        \item \textbf{Hardware Specificity.}  
              Results reflect performance on NVIDIA GPUs.  
              Latency and efficiency trends may differ on CPU-only or  
              low-power edge devices where memory patterns dominate.

        \item \textbf{Intrinsic Architectural Constraints.}  
              Some architectures carry structural limitations that  
              hinder convergence under standardized training settings.  
              Examples include SqueezeNet's aggressive compression,  
              DenseNet's feature-map accumulation, and the depth of  
              large ResNet variants. These models may require  
              architecture-specific schedules to reach full potential.
    \end{itemize}
\end{frame}

% Final Slide
\begin{frame}
  \frametitle{Thank You!}
  \InfContacts
\end{frame}

\end{document}